\problemname{Lördagsgodis}
Enhörningen Omar älskar godis! Han har precis fått $P$ kronor av sina mamma för att köpa lördagsgodis.
Eftersom Omar är smart kommer han att försöka spendera pengarna så bra som möjligt. I butiken han brukar
handla godis i finns det $N$ stycken godispåsar. Varje godispåse har ett visst pris och en viss kategori.
Eftersom Omar gillar godis väldigt mycket så tycker han \textit{nästan} att alla godispåsar är lika goda. 
Han har nämligen ingen preferens för vilken kategori av godis han ska köpa, men han tycker att det blir 
lite för enformigt om han köper mer än $K$ påsar av samma kategori. 

Hjälp Omar att köpa så många godispåsar som möjligt, vars totala kostnad är mindre än eller lika med $P$.
Se dessutom till att det inte finns mer än $K$ påsar av samma kategori bland de han köpte. Du behöver
bara skriva ut hur många påsar han kan köpa, inte vilka.

\section*{Indata}
Den första raden i indatan innehåller heltalet $N$ ($1 \leq N \leq 10^5$), antalet godispåsar som Omar
kan välja mellan. 

Nästa rad innehåller heltalet $P$ ($1 \leq P \leq 10^8$), antalet kronor som Omar har i godisbudget.

Nästa rad innehåller heltalet $K$ ($1 \leq K \leq N$), antalet godistyper som Omar kan tänka sig
äta av varje kategori.

På nästa rad står heltalen $c_1, c_2, \dots, c_N$ ($1 \leq c_i \leq P$), där $c_i$ betyder att
godispåse $i$ kostar $c_i$ kronor.

På sista raden står heltalen $t_1, t_2, \dots, t_N$ ($1 \leq t_i \leq 10^5$). Det finns exakt $10^5$
kategorier av godispåsar, och vi benämner varje med ett heltal. Varje $t_i$ betyder att
godispåse $i$ tillhör kategori $t_i$.

\section*{Utdata}
Skriv ut ett heltal: antalet godispåsar som Omar kan köpa om han använder sin budget optimalt.

\section*{Poängsättning}
Din lösning kommer att testas på en mängd testfallsgrupper.
För att få poäng för en grupp så måste du klara alla testfall i gruppen.

\noindent
\begin{tabular}{| l | l | p{12cm} |}
  \hline
  \textbf{Grupp} & \textbf{Poäng} & \textbf{Gränser} \\ \hline
  $1$    & $20$       & Affisch-sample. \\ \hline
  $2$    & $5$        & $K=N$ och du har råd att köpa allt godis. \\ \hline
  $3$    & $15$       & Du har råd att köpa allt godis. \\ \hline
  $4$    & $5$        & $c_i=1$, det vill säga kostar varje sorts godis $1$ krona. \\ \hline
  $5$    & $17$       & $K=N, N \leq 500$. \\ \hline
  $6$    & $21$       & $K=N$ \\ \hline
  $7$    & $9$        & $N \leq 500$. \\ \hline
  $8$    & $8$        & Inga ytterligare begränsningar. \\ \hline
\end{tabular}

\section*{Förklaring av exempelfall 1}
I exempelfall 1 går det jättesnabbt att lösa problem, men tar jättelång tid att skapa nya. Optimala strategin
är därmed att bara lösa problem. För att lösa mer än hälften måste han lösa $51$ stycken. Eftersom $A=51$ blir svaret
$51$.

I exempelfall 2 måste 2 problem lösas för att lösa mer än hälften.

I exempelfall 3 är den optimala strategin att lösa 11 problem och skapa 1 problem. Detta tar sammanlagt $11 \cdot 7 + 1 \cdot 6=83$ sekunder.
Det finns ingen annan strategi som är snabbare.
